% ===============================================================
%  Chapter 4 -- How Cache Coherence Is Maintained        (Part I)
% ===============================================================
\chapter{How Cache Coherence Is Maintained}
\label{ch:coherence-basics}

\begin{figure}[htbp]
\centering
\adjustbox{max width=\linewidth}{%
\begin{forest}
  chapter map,
  [Chapter 4
    [What coherence means [Completion {\&} order axioms] [Not consistency]]
    [Block states [Event / action]]
    [Snooping and directories [Bus orders messages] [Sharer list]]
    [Update vs invalidate [Decision tree]]
  ]
\end{forest}}
\caption[]{Structure of Chapter~4: what a coherence
protocol must guarantee, then the parts every protocol is built from.}
\label{fig:ch4-map}
\end{figure}

% ---------------------------------------------------------------
\section{What Coherence Means}
\label{sec:coh-meaning}
\label{sec:coh-order}
\label{sec:coh-consistency}
\sectionmap{[Coherence [Concerns one address] [Completion: no write is lost]
  [Order: one order for all threads] [{Consistency: a separate question}]]}
\index{coherence axioms}\index{memory consistency}\index{Sarangi textbook}%

To a program, a set of private caches has to look like one cache. Sarangi states
the requirement in terms of the accesses to a single
address~\cite[\S12.4.2, Def.~122]{sarangi2025basic}\index{coherence axioms}:

\begin{keybox}[The two axioms of coherence]
\emph{Completion}: a write must eventually complete, and the memory system may
never lose it.\\[3pt]
\emph{Order}: all threads see the writes to one address in the same order.
\end{keybox}

However many copies of a location the machine holds, it must behave like a single
location whose value changes in one agreed sequence. \Cref{fig:order-example}
draws Sarangi's example of the order axiom~\cite[\S12.4.2, Ex.~151]{sarangi2025basic}.
Thread~1 writes \texttt{x=1} and then \texttt{x=2}, and thread~2 reads \texttt{x}
twice. Every pair of results can occur except $(2,1)$, which would mean thread~2
saw the second write before the first. A coherence protocol is the hardware that
enforces both axioms for every block in every cache.

Coherence says nothing about writes to different addresses. Whether another
thread may see a new value of \texttt{y} together with an old value of \texttt{x}
after thread~1 wrote \texttt{x} and then \texttt{y} is a question of memory
consistency~\cite[\S12.4.3]{sarangi2025basic}. We keep the consistency model fixed
and study coherence only.

% ---------------------------------------------------------------
\section{Block States and State Machines}
\label{sec:coh-fsm}
\sectionmap{[States [Each cached block has a state] [Events cause transitions]
  [Written as event / action]]}

All the protocols in this report work the same way. Every block in every cache
has a small state stored next to its tag, and the cache controller runs a finite
state machine for each block. When an event arrives, the controller performs an
action and may move the block to a new state~\cite[\S12.4.6]{sarangi2025basic}.
Events come either from the cache's own processor (read, write, evict) or from
other caches through the interconnect (\cref{fig:fsm-notation}). A block that is
not in the cache counts as Invalid.

\begin{figure}[htbp]
\centering
\begin{minipage}[c]{0.5\linewidth}\centering
\adjustbox{max width=\linewidth}{%
\begin{tikzpicture}[font=\sffamily\small]
  \node[font=\sffamily\bfseries, text=tealC!80!black] at (0,2.0) {Thread 1};
  \node[font=\sffamily\bfseries, text=plumC] at (5.6,2.0) {Thread 2};
  \node[draw=tealC, fill=tealC!10, rounded corners=2pt] at (0,1.35) {\texttt{x = 1}};
  \node[draw=tealC, fill=tealC!10, rounded corners=2pt] at (0,0.7) {\texttt{x = 2}};
  \node[draw=plumC, fill=plumC!10, rounded corners=2pt] at (5.6,1.35) {\texttt{t1 = x}};
  \node[draw=plumC, fill=plumC!10, rounded corners=2pt] at (5.6,0.7) {\texttt{t2 = x}};
  \draw[-{Stealth}, thick, amberC] (0.8,1.35) to[out=-15,in=195] (4.65,1.35);
  \draw[-{Stealth}, thick, amberC] (0.8,0.7) to[out=-15,in=195] (4.65,0.7);
  \draw[thick, coreGrey!40] (-1.0,0.15) -- (6.6,0.15);
  \node[text=leafC!80!black, align=center] at (2.8,-0.5) {Allowed: $(0,0)$, $(0,1)$, $(0,2)$,\\$(1,1)$, $(1,2)$, $(2,2)$};
  \node[text=brickC] at (2.8,-1.25) {Forbidden: $(2,1)$, writes seen out of order};
\end{tikzpicture}}
\caption[The order axiom]{The order axiom: only $(2,1)$ is impossible.}
\label{fig:order-example}
\end{minipage}\hfill
\begin{minipage}[c]{0.45\linewidth}\centering
\adjustbox{max width=\linewidth}{%
\begin{tikzpicture}[font=\sffamily\small, >=Stealth,
    st/.style={draw, thick, circle, minimum size=1.2cm, font=\sffamily\bfseries}]
  \node[st, fill=skyC!16, draw=skyC] (A) at (0,0) {A};
  \node[st, fill=leafC!16, draw=leafC] (B) at (5,0) {B};
  \draw[->, thick] (A) -- node[above]{\textit{event} / \textit{action}} (B);
  \node[draw, thick, rounded corners=2pt, fill=slateC!12, draw=slateC, align=center, font=\sffamily\scriptsize]
       (p) at (-0.2,-2.0) {from this cache's processor:\\read, write, evict};
  \node[draw, thick, rounded corners=2pt, fill=plumC!12, draw=plumC, align=center, font=\sffamily\scriptsize]
       (o) at (5.2,-2.0) {from other caches:\\read miss, write miss};
  \draw[->, thick, slateC] (p) -- (2.3,-0.25);
  \draw[->, thick, plumC] (o) -- (2.7,-0.25);
\end{tikzpicture}}
\caption[State-machine notation]{Notation for every protocol in this report.}
\label{fig:fsm-notation}
\end{minipage}
\end{figure}

% ---------------------------------------------------------------
\section{Snooping and Directory Protocols}
\label{sec:coh-snoop}
\label{sec:coh-dir}
\sectionmap{[Organisation [Snooping: every cache watches the bus]
  [Bus order gives the order axiom] [Directory: sharer list per block]]}
\index{snooping protocol}\index{directory protocol}\index{coherence axioms}%

In a \gls{snoop}\index{snooping protocol} all private caches sit on one shared bus.
One cache puts a message on the bus at a time and every other cache reads
it~\cite[\S12.4.6]{sarangi2025basic}. Each controller watches (snoops) the bus and
reacts to messages about blocks it holds. An arbiter grants the bus to requesters
in FIFO order~\cite[\S13.4.3]{sarangi2025basic} (\cref{fig:snoop-bus}). Since only
one cache transmits at a time, all caches see all messages in the same order, and
the order axiom holds automatically~\cite[\S12.4.6]{sarangi2025basic}. One
consequence is easy to miss. A cache that wants to write must wait for its own
request to appear on the bus before updating its copy. If it wrote first, two
caches writing at the same moment could each apply their own write first and see
the two writes in opposite orders~\cite[\S12.4.6, Question~8]{sarangi2025basic}.

\begin{figure}[htbp]
\centering
\begin{minipage}[c]{0.55\linewidth}\centering
\adjustbox{max width=\linewidth}{%
\begin{tikzpicture}[font=\sffamily\scriptsize,
    cpu/.style={draw, thick, rounded corners=2pt, fill=slateC!14, draw=slateC,
                minimum width=1.4cm, minimum height=0.55cm},
    cache/.style={draw, thick, rounded corners=2pt, minimum width=1.9cm,
                  minimum height=0.5cm}]
  \foreach \i/\c in {0/tealC,1/plumC,2/amberC,3/leafC} {
    \node[cpu] (c\i) at (\i*2.4,2.3) {P\i};
    \node[cache, fill=\c!15, draw=\c] (h\i) at (\i*2.4,1.5) {L1 + snooper};
    \draw[thick] (c\i) -- (h\i); \draw[thick] (h\i) -- (\i*2.4,0.8);
  }
  \draw[line width=2pt, slateC] (-1.0,0.8) -- (8.2,0.8);
  \node[text=slateC, anchor=south east] at (8.2,0.85) {shared bus};
  \node[draw, thick, rounded corners=2pt, fill=brickC!12, draw=brickC,
        minimum width=2.2cm, align=center] (arb) at (1.2,-0.45) {Arbiter\\FIFO queue};
  \draw[thick, brickC, <->] (arb) -- (1.2,0.8);
  \node[draw, thick, rounded corners=2pt, fill=leafC!14, draw=leafC,
        minimum width=2.4cm] (mem) at (6.0,-0.45) {Lower level / memory};
  \draw[thick] (mem) -- (6.0,0.8);
\end{tikzpicture}}
\caption[A snooping system]{Four private caches on a shared bus. The arbiter
lets one cache transmit at a time.}
\label{fig:snoop-bus}
\end{minipage}\hfill
\begin{minipage}[c]{0.42\linewidth}\centering
\adjustbox{max width=\linewidth}{%
\begin{tikzpicture}[font=\sffamily\scriptsize,
    cache/.style={draw, thick, rounded corners=2pt, minimum width=1.3cm, minimum height=0.5cm}]
  \foreach \i/\c in {0/tealC,1/slateC,2/plumC,3/plumC} { \node[cache, fill=\c!14, draw=\c] (h\i) at (\i*1.7,2.2) {cache \i}; }
  \node[draw, thick, rounded corners=2pt, fill=amberC!16, draw=amberC!90!black,
        minimum width=3.0cm, minimum height=0.8cm] (d) at (2.55,0.4) {Directory};
  \draw[-{Stealth}, thick, tealC] (h0) -- node[left, pos=0.4, text=tealC!80!black]{request} (d);
  \draw[-{Stealth}, thick, plumC] (d) -- (h2);
  \draw[-{Stealth}, thick, plumC] (d) -- node[right, pos=0.6, text=plumC]{forward} (h3);
  \draw[dashed, coreGrey!50] (d) -- (h1);
  \begin{scope}[xshift=0.3cm, yshift=-1.3cm]
    \draw[thick, fill=slateC!10] (0,0) rectangle (1.4,0.6); \node at (0.7,0.3) {block addr.};
    \draw[thick, fill=amberC!18] (1.4,0) rectangle (2.2,0.6); \node at (1.8,0.3) {state};
    \foreach \b/\v in {0/0,1/0,2/1,3/1} {
      \ifnum\v=1 \def\fc{plumC!30}\else\def\fc{white}\fi
      \draw[thick, fill=\fc] (2.2+\b*0.45,0) rectangle +(0.45,0.6);
      \node at (2.425+\b*0.45,0.3) {\v};
    }
    \node[text=coreGrey] at (3.1,-0.3) {sharer bit per cache};
  \end{scope}
\end{tikzpicture}}
\caption[A directory protocol]{The directory forwards cache~0's request only to
the recorded sharers, caches~2 and~3.}
\label{fig:directory}
\end{minipage}
\end{figure}

Broadcasting every request wastes bandwidth when only one or two caches hold the
block. A \gls{directory}\index{directory protocol} keeps a record for each block of
the caches that may hold a copy, usually one bit per
cache~\cite[\S12.4.6]{sarangi2025basic}. Requests go to the directory, which
forwards them only to the caches concerned (\cref{fig:directory}). Censier and
Feautrier proposed the scheme in 1978~\cite{censier1978directory}.
\Cref{tab:snoop-dir} sums up the trade-off.

\begin{table}[htbp]
\centering
\caption{Snooping and directory organisations compared.}
\label{tab:snoop-dir}
\small
\begin{tabular}{@{}l p{4.8cm} p{4.8cm}@{}}
\toprule
 & Snooping & Directory \\
\midrule
Who sees a request & every cache & only the recorded sharers \\
Ordering of messages & provided by the bus & enforced by the directory \\
Extra storage & none & a directory entry per block \\
Interconnect & shared bus & any (ring, mesh, \ldots) \\
Scales to & a small number of cores & many cores \\
\bottomrule
\end{tabular}
\end{table}

% ---------------------------------------------------------------
\section{Write-Update and Write-Invalidate}
\label{sec:coh-update-invalidate}
\label{sec:coh-choose}
\sectionmap{[Two strategies [Update: send new value to all copies]
  [Invalidate: destroy other copies] [MSI family is write-invalidate]]}
\index{MSI}\index{snooping protocol}\index{directory protocol}\index{write-update}\index{write-invalidate}\index{Dragon protocol}\index{Firefly protocol}\index{Sarangi textbook}%

A protocol must also decide what happens to the other copies when one cache writes
a block~\cite[\S12.4.6]{sarangi2025basic}. Under \gls{update}\index{write-update}
the writer sends each new value out and every other copy is updated in place, so
each write to shared data costs a message. Under
\gls{invalidate}\index{write-invalidate} the writer first removes every other
copy, after which it can write as often as it likes without further messages. At
any moment a block then has either one writer or any number of readers. Sarangi
justifies write-invalidate by noting that most blocks are not shared at the moment
they are written. This choice is separate from the snooping-or-directory choice.
The MSI family of Chapter~5 is write-invalidate and exists in both snooping and
directory forms, while Dragon and Firefly (Part~III) are snooping write-update
protocols. \Cref{fig:choose-tree} puts both choices into one decision tree. Real
designs mix them, for example by snooping within a cluster and using a directory
between clusters. From here on every protocol is write-invalidate unless stated
otherwise.

\begin{figure}[htbp]
\centering
\adjustbox{max width=0.8\linewidth}{%
\begin{forest}
  decision tree,
  [Few cores sharing\\one bus?
    [Snooping, el=yes
      [Same core usually\\rewrites before\\others read?
        [{Snooping +\\invalidate\\(MSI family)}, el=yes]
        [{Snooping +\\update}, el=no]]]
    [Directory, el=no
      [Same core usually\\rewrites before\\others read?
        [{Directory +\\invalidate}, el=yes]
        [{Directory +\\update}, el=no]]]]
\end{forest}}
\caption[Decision tree: choosing a coherence organisation]{The two design
questions of this chapter. Real programs usually answer the second question
``yes'', which is why write-invalidate is the common choice.}
\label{fig:choose-tree}
\end{figure}
